\documentclass{article}
\usepackage[utf8]{inputenc}
\usepackage[T1]{fontenc}
\usepackage[margin=2cm]{geometry}
\usepackage{hyperref}
\sloppy
\begin{document}
\section*{Historical source}
Encyclopaedia Britannica scan, source row 2402. Parse 3.5 extracted the text. Super 3.5 VL corrected it from an enlarged crop of each block. It can still contain recognition errors.
Page headers, footers, and page numbers are removed.

\subsection*{Untitled continuation (E0215)}
\textit{Continues an entry from an earlier page.}

Comparative Anatomy.

Sacral ver-
tebræ.

Coccygeal
or caudal
vertebre.

Occuges
or catalans
vertebre.

The number of sacral vertebrae is still more various, even in the species of the same genus. Thus, while in several of the ape genus, in the lori, in the vampyre bat, the colugo (galeopithecus), the coati, and two of the didelphis, there is one sacral vertebra only, most of the apes have sacra consisting of 2, 3, 4, 5, or 6 pieces; the majority of other animals have 3 sacral vertebrae; the hedgehog, porcupine, guinea pig, paca, hare, tiger, several of the murine genus, the ant-eater, rhinoceros, camel, dromedary, chamois, goat, sheep, and ox, have 4; the elephant has 5; the Ai 6; the Unau 7; and in the mole, white bear, and quagga, they also amount to 7. The frequency of the three sacral vertebrae in the lower animals shows that Galen, who ascribes only 3 to the human subject, must have derived this inference from the former.

These vertebrae are in the mammals narrower than in
man, and their direction forms with the spine, instead of
receding backwards, a straight line; an arrangement evi-
dently connected with the horizontal position of the for-
mer. The shape of the sacrum in the lower mammals is
that of an elongated triangle; and it is further remarkable,
that in those species which occasionally assume the erect
attitude on the hind leg, as apes, bears, and sloths, the
width of the sacrum is proportionally greater. The sa-
cral spines, which are short in man and the ape, become
longer in the Zoophaga, and form a continuous ridge in the
rhinoceros, most ruminants, and especially in the mole. In
the vampire bat the sacrum forms a long compressed
cone, the extremity of which is united to the ischial tu-
berosities, without sustaining a coccyx. The seal has two
sacral bones; but the Cetacea, e. g. the dolphin and por-
poise, are void both of sacrum and coccyx.

The coccygeal bones constitute the tail of the lower animals, and in many instances they are extremely numerous. The smallest number is 3, which is that of the magot (simia sylvanus, pithecus, et inus) or Barbary ape; and the greatest yet known is that of the ant-eater, in which they amount to 40, and the long-tailed manis, in which they amount to 45. Next to these may be placed that of the coaita 32, the baboon 31, the phalanger (delphis orientalis) 30, the marmoset (delphis murina) 29, the pangolin 28, the silky monkey (simia rosalia) and black rat 26, the weeping monkey and howling ape 25; the panther, mouse, dormouse, and elephant, 24; the lion, beaver, water-rat, Norway rat, and field-rat, 23; the flying-cat, puma, cat, dog, marmot, and rhinoceros, 22; the otter, 21; the Chinese monkey, raccoon, civet, hare, and rabbit, 20; the tiger and wolf, 19; the macauco, glutton, marten, fat dormouse, dromedary, giraffe, and quagga, 18; the tarsier, shrew, camel, and horse, 17; and other genera and species, without any determinate order, descending so low as to 9, 8, 7, 6, and 4. The quilled duckbill (echidna, ornithorhynchus hystrix) has only 12 caudal vertebrae, while the common one (ornithorhynchus paradoxus) has at least 20. The gibbon and vampyre bat are the only mammiferous animals, excepting the Cetacea, in which there are no coccygeal bones. It sometimes happens that a monkey or opossum loses a portion of its tail, when the truncated end is converted into a knotty excrescence, sometimes carious, always different from the taper point of the last coccygeal vertebra; and in this case it is difficult to determine the exact species.

In the Cetacea, in which the absence of pelvis affords
no mark to distinguish the lower vertebrae into lumbar,
sacral, and coccygeal, those below the dorsal may be re-

regarded as lumbo-coccygeal; and their number is estimated
by deducting that of the cervical and costal from the total
number. The following table, which shows the number
of the costal, the lumbo-coccygeal, and the total number
of vertebrae, indicates that their number varies much in
various genera of this family.

various genera of this family.
Lamantin.....16 24 46
Dugong.....16 28 46
Dolphin.....13 47 67
Turso.....13 38 58
D. Globiceps.....11 37 56
D. Griseus.....12 42 61
D. Gangeticus.....11 28 46
Porpoise.....13 40 60
Narwal.....12 35 54
Hyperoodon.....9 29 45
Cachalot or White
Whale.....14-15 39 60
Greenland Whale.....15 37 69
Rorqual.....14 31 52

In general, however, if we reckon the first 2, 3, or 5 verte-
bre after the costal as lumbar, it may be said that the
caudal vertebrae of the Cetacea vary from 22 or 25, which
are the numbers respectively in the lamantin and dugong,
to 34, 38, and 42, at which they may be estimated in the
dolphin. We shall see that, in the dugong at least, we are
guided in this estimate by the rudimentary bones of the
pelvis.

The coccygeal or caudal vertebrae of the Mammalia may be distinguished into two kinds; those which contain a canal in continuity with that of the vertebral column and sacrum, and those in which the pieces are solid. The first, which are next the sacrum, have articular, transverse, and spinous processes, distinct in proportion as the animals move their tails. The latter are generally prismatic in shape, diminish in size towards the extremity, and have only slight tubercles for muscular attachments. Animals with prehensile tails, as the American ape (sapajous), have above, at the base of the body of each vertebra, two small tubercles, between which pass the tendons of the flexor muscles. By means of this mechanism these animals can twine the tail round the branch of a tree with sufficient force to support the weight of the body.

The MAMMALIA with long mobile tails have often two
or three small supernumerary bones placed on the lower
surface of the junctions of several of the coccygeal ver-
tebræ, from the 3d or 4th to the 7th or 8th. These
sesamoid bones give attachment to muscles. In the
beaver, which employs its tail as a trowel, the transverse
processes are remarkable for size, while the lower spinous
processes are larger than the upper ones,—an arrangement
which enables it to depress the tail forcibly when it beats
the ground.

The shape of the chest in the Mammalia varies in general as the clavicles are present or wanting. In animals provided with clavicles, as the Quadrumana, bats, the squirrel, beaver, mole, ant-eater, hedgehog, and sloth, the shape of the chest approaches to the human, or is conical, and flattened before and behind. In those void of clavicles it is compressed laterally, from the smaller incurvation of the ribs; and the sternum makes a remarkable prominence, so that the transverse or intercostal diameter is less proportionally, and the sterno-vertebral is greater proportionally, than in man and the claviculated animals. In the long-legged animals, as the giraffe and those of the stag kind, this prominence of the sternum is sufficient to give it a keel-like appearance (thorax carinatus). In the carnivorous animals the chest presents its greatest longitudinal extent.

The number and shape of the ribs varies in the different tribes. In number, indeed, the ribs always correspond with that of the costal vertebrae. Thus, in the Quadrumana, Zoophaga, Rodentia, Edentata, and Ruminantia, they vary from 12 to 15 pair, with only three exceptions, the glutton, hyena, and ant-eater. In the Chinese monkey, common bat, and armadillo, they are a pair less than in


\section*{Source record}
\begin{description}
\item[Scan ID] \texttt{britannica:309f765b1ecc83e0a7f7762e44b4a927a7555443:train:2402}
\item[Image file] \texttt{page-2402-5030b31ea901.jpg}
\item[Image SHA-256] \texttt{5030b31ea90173a447b4749960488cdeab085e4ff2f4bf49cf473614e50ce29e}
\item[Dataset] NationalLibraryOfScotland/encyclopaedia\_britannica\_illustrated at revision \texttt{309f765b1ecc83e0a7f7762e44b4a927a7555443} (CC0-1.0)
\item[Extraction] model \texttt{nvidia/nemotron-3.5-parse}, raw reply \texttt{extract/raw/page-2402-5030b31ea901.json}
\item[Raw reply SHA-256] \texttt{9f34d69e9ccfd986114c9a75010547bd871ce6d804c770fd86f437aaab292fae}
\item[Correction] model \texttt{nvidia/Nemotron-3.5-Super-VL}, block crops, raw replies \texttt{correct/raw/crop\_nothink/page-2402-5030b31ea901-b*.json}; parse\_fallback: 1, super\_vl: 13
\item[Running header] A N A T O M Y. 3
\end{description}
\end{document}
